% !TeX root = ../main.tex
\chapter{Discussion}
\label{chap:discussion}

% I think "forces and design trade-off" should be in discussion part
\section{Forces and Design Trade-offs}
\label{sec:cd4ai-forces}

\noindent\emph{Working formulation: the following forces will be revisited
as the case studies are developed.}

Testability, robustness, and evolvability are desired qualities that can
reinforce one another. Testability helps teams assess whether a proposed change
preserves expected behavior before deployment. Robustness concerns the
application's ability to maintain acceptable behavior under varying conditions,
while evolvability concerns its ability to accommodate change. The evaluation
suite must also evolve to remain useful as the application changes.

These qualities need not oppose one another in every pairwise comparison.
The forces considered here instead express tensions between desired outcomes
and the constraints on achieving them. The pattern aims to address three
such tensions:

\begin{itemize}
  \item \textbf{Confidence in changes and delivery speed:} teams need evidence
    that changes preserve expected behavior, but running more extensive
    evaluations requires execution time and maintenance effort. The validation
    process must fit the time and resources available for frequent delivery.
  \item \textbf{Failure discovery and review capacity:} teams need to discover
    failures beyond those covered by existing tests, but collecting more
    candidate interactions increases the curation workload. Broad discovery
    must be balanced with the capacity to review candidates and prepare useful
    evaluation cases.
  \item \textbf{Behavior preservation and adaptation:} historical regression
    cases help check established behavior, but changing requirements can make
    some expectations outdated. Teams must distinguish unintended regressions
    from intended changes and maintain the suite accordingly.
\end{itemize}

Evaluation, monitoring, and curation work together to address these tensions.
Evaluation supplies evidence about proposed changes, monitoring supplies
interactions for investigation, and curation determines which interactions
should become evaluation cases and what behavior those cases should check.
The cycle does not remove the costs or trade-offs; its implementation must
establish how these responsibilities fit the team's delivery process.

% Writing notes:
% Keep the desired qualities distinct from the concrete tensions above.
% Do not claim that greater testability necessarily lowers robustness or
% evolvability. Better testing can support both.
% Validate this working formulation against Chapter 5: look for evidence of
% evaluation cost, candidate selection, review workload, and changes to expected
% behavior. Do not present these as findings before examining the evidence.
% Explain how the stages address the tensions without forcing a one-to-one
% correspondence between a stage and a force. Connected paragraphs are enough;
% pairwise-comparison subsections are not necessary.
% Pattern-writing guidance discussed during review:
% https://www.hillside.net/index.php/a-pattern-language-for-pattern-writing
% Add a formal citation if this guidance is used in the final argument.


% I think "consequences and limitations" should be in discussion part
\section{Consequences and Limitations}
\label{sec:cd4ai-consequences}

\begin{itemize}
  \item Discuss the intended benefit of an evaluation suite informed by actual
    use, separating expected benefits from measured outcomes.
  \item Discuss the costs of observation, review, and maintaining and running
    the growing suite.
  \item Explain limits arising from unobserved interactions, missed signals,
    reviewer errors, and variable behavior. Avoid guarantees of complete
    coverage or prevention of every recurrence.
  \item Introduce the case studies as an examination of how these responsibilities
    and connections appear in different systems and workflows.
\end{itemize}

\textbf{Draft outline.} Compare the cases, discuss the research question,
and explain limitations and threats to validity.

